\documentclass[11pt,letterpaper]{article}
\def\SuiteDocTitle{User Stories and Agile Requirements Engineering}
\def\SuiteDocCode{TR-01}
% Shared professional style: pdfLaTeX compatible, no external font files.
\usepackage[T1]{fontenc}
\usepackage[utf8]{inputenc}
\usepackage{lmodern}
\usepackage{microtype}
\usepackage[margin=0.83in,headheight=16pt,footskip=28pt]{geometry}
\usepackage[table]{xcolor}
\definecolor{navy}{HTML}{12304A}
\definecolor{teal}{HTML}{087F83}
\definecolor{gold}{HTML}{B78633}
\definecolor{blueback}{HTML}{F0F5F8}
\definecolor{tealback}{HTML}{EEF8F6}
\definecolor{amberback}{HTML}{FFF7E8}
\definecolor{grayback}{HTML}{F6F7F8}
\usepackage{array,tabularx,longtable,booktabs}
\usepackage{enumitem}
\usepackage{amssymb}
\usepackage{titlesec}
\usepackage{fancyhdr,lastpage}
\usepackage{tcolorbox}
\tcbuselibrary{breakable}
\usepackage{pgffor}
\usepackage{hyperref}
\hypersetup{colorlinks=true,linkcolor=teal,urlcolor=teal,pdfauthor={Engineering Learning Series},pdftitle={\SuiteDocTitle}}
\setlength{\parindent}{0pt}
\setlength{\parskip}{5.5pt}
\setlist{topsep=3pt}
\setcounter{tocdepth}{2}
\titleformat{\section}{\Large\bfseries\color{navy}}{\thesection}{0.7em}{}
\titleformat{\subsection}{\large\bfseries\color{teal}}{\thesubsection}{0.6em}{}
\titleformat{\subsubsection}{\normalsize\bfseries\color{navy}}{\thesubsubsection}{0.6em}{}
\pagestyle{fancy}
\fancyhf{}
\fancyhead[L]{\scriptsize\color{navy}AGILE REQUIREMENTS ENGINEERING}
\fancyhead[R]{\scriptsize\color{teal}\SuiteDocCode}
\fancyfoot[L]{\scriptsize Independent learning and implementation materials}
\fancyfoot[R]{\scriptsize Page \thepage\ of \pageref{LastPage}}
\renewcommand{\headrulewidth}{0.35pt}
\renewcommand{\footrulewidth}{0pt}
\newcommand{\SuiteHeader}[2]{%
{\color{teal}\small\bfseries ENGINEERING DOCUMENTATION SUITE \hfill #2}\par\vspace{8pt}
{\color{navy}\LARGE\bfseries\SuiteDocTitle}\par\vspace{5pt}
{\color{black!68}\normalsize #1}\par\vspace{8pt}
{\color{gold}\rule{\linewidth}{2pt}}\par\vspace{8pt}}
\newcommand{\WritingLines}[1]{\par\foreach \n in {1,...,#1}{\noindent\rule{\linewidth}{0.35pt}\par\vspace{0.28cm}}}
\newcommand{\FreeSpace}[1]{\begin{tcolorbox}[colback=white,colframe=black!25,height=#1,boxrule=0.5pt]\end{tcolorbox}}
\newcommand{\Field}[1]{\noindent\textbf{#1}\quad\rule{0.42\linewidth}{0.35pt}\par\vspace{0.3cm}}
\newcommand{\CheckItem}[1]{\noindent$\square$\quad #1\par\vspace{2pt}}
\newcommand{\SmallNote}[1]{{\footnotesize\color{black!62}#1}\par}
\newcommand{\Guidance}[1]{\begin{tcolorbox}[colback=blueback,colframe=teal,breakable,title=Guidance]#1\end{tcolorbox}}
\sloppy

\begin{document}
\SuiteHeader{Technical reference | lifecycle, heuristics, verification and governance}{Reference}
\tableofcontents\medskip\noindent\rule{\linewidth}{0.4pt}\medskip
\section{Purpose, provenance and boundaries}
This reference converts the principles and examples of Mike Cohn’s User Stories Applied (2004) into practical, original guidance. Concepts attributed to the book include Card–Conversation–Confirmation, INVEST, user role modeling, story workshops, acceptance tests, relative estimation, release/iteration planning, story smells, and the extended case-study workflow. Modern extensions in this suite include explicit evidence lineage, repository-managed templates, secure delivery checks, and automated acceptance-test traceability; these additions are not claims about the 2004 text.

\section{The requirements artifact system}
\begin{longtable}{|p{0.20\linewidth}|p{0.29\linewidth}|p{0.37\linewidth}|}\hline
\textbf{Artifact} & \textbf{Primary question} & \textbf{Minimum durable information} \\ \hline\endfirsthead\hline
\textbf{Artifact} & \textbf{Primary question} & \textbf{Minimum durable information} \\ \hline\endhead
Outcome / product goal & Why should anyone invest? & Target behavior, beneficiaries, success measure, constraints \\ \hline
User role model & Who needs what? & Role, relevant goal, context, differentiated permissions \\ \hline
Story / product backlog item & What thin outcome can be delivered next? & Role, capability, value, status, conversation needed \\ \hline
Conversation decision & What has been agreed? & Decision, rationale, alternatives if meaningful, owner, date \\ \hline
Acceptance example & How will we know? & Preconditions, action, observable outcome, test oracle \\ \hline
Constraint / quality attribute & What must be true across features? & Scope, measurable threshold, environment, owner, method \\ \hline
Use case or process model & What end-to-end behavior must endure? & Trigger, actors, main/alternate flows, boundaries \\ \hline
ADR or interface contract & Why this technical arrangement? & Problem, decision, consequences, interfaces, version \\ \hline
Release evidence & What was actually accepted? & Commit/build, tests, results, exceptions, approvals \\ \hline
\end{longtable}
\section{Card, Conversation, Confirmation}
Card: a short prompt to plan a possible outcome. Conversation: the ongoing collaboration that resolves roles, business rules, error cases and tradeoffs. Confirmation: concrete tests and examples establishing whether the agreed behavior is present. The approach deliberately avoids treating the initial story sentence as exhaustive requirements documentation.

\begin{tcolorbox}[colback=tealback,colframe=navy,title={Original example: service request},breakable]
Card — As a requester, I want to submit a service request so that I receive a tracking reference. Conversation — Can drafts be saved? What information is mandatory? Can users upload evidence? Confirmation — A valid request returns a unique reference and is visible in the requester’s history; an invalid request shows actionable validation errors.

\end{tcolorbox}
\section{Story quality and decomposition}
\begin{itemize}[leftmargin=*,itemsep=2pt]
\item Independent: can be ordered and developed with few unnecessary couplings; document true external dependencies.
\item Negotiable: prompts discovery rather than fixing implementation prematurely.
\item Valuable: a named role or purchaser obtains a discernible benefit.
\item Estimatable: unknown scope and risks have been discussed or explicitly investigated.
\item Small: near-term items fit a short iteration and yield a demonstrable result.
\item Testable: expected behavior has objective and discriminating evidence.
\end{itemize}
For vertical slicing, partition a large capability by business rule, workflow path, data variation, role, or minimal-to-richer behavior. Each slice should support an end-to-end outcome. Avoid technical-layer slices that leave no functioning feature. Keep optional capabilities separate from core behavior.

\section{Role modeling and elicitation}
Start with roles defined by behavior rather than organization chart. Collect actual tasks and situational context by interviews, observation, questionnaires and story-writing workshops. Distinguish first-hand evidence from proxy statements; record stakeholder authority separately from empirical user knowledge. A user-role card should include goals, constraints and an evidence source.

\section{Acceptance engineering}
Acceptance examples are the shared specification of behavior near implementation. The product/customer side specifies intended outcomes and participates in review; testers and developers collaborate on specificity, coverage and automation. Include normal flow, invalid input, denied authorization, boundaries, state changes, and error recovery when relevant. Pair these examples with developer-owned unit/integration/security tests; do not confuse one level with another.

\begin{longtable}{|p{0.20\linewidth}|p{0.29\linewidth}|p{0.37\linewidth}|}\hline
\textbf{Layer} & \textbf{Example check} & \textbf{Primary owner} \\ \hline\endfirsthead\hline
\textbf{Layer} & \textbf{Example check} & \textbf{Primary owner} \\ \hline\endhead
Story acceptance & Requester receives tracking reference after valid submission & Product representative + tester \\ \hline
API contract & Response status and schema match published interface & Engineering \\ \hline
Unit / integration & Validation and persistence have correct edge behavior & Engineering \\ \hline
Security / privacy & Requester cannot read another requester’s record & Security + engineering \\ \hline
Quality constraint & 95th percentile submission latency is within agreed limit & Engineering + product/service owner \\ \hline
Operational readiness & Failed delivery is observable and recoverable & Operations + engineering \\ \hline
\end{longtable}
\section{Estimating and planning operating model}
\begin{enumerate}[leftmargin=*,itemsep=3pt]
\item Define a clear release goal and identify the smallest complete user journey that serves it.
\item Create an ordered backlog based on user value, risk, dependencies and enablement work.
\item Establish a team-local reference scale, estimate collectively and surface disagreements.
\item Use observed throughput or completed points as a forecast input, explicitly showing uncertainty.
\item During iteration planning, clarify acceptance examples, break stories into tasks and confirm capacity.
\item Inspect accepted outcomes at review, distinguish scope change from estimate change, and replan.
\end{enumerate}
\begin{tcolorbox}[colback=tealback,colframe=navy,title={Illustrative forecast (not a commitment)},breakable]
If a team completed 13, 15 and 12 points in three comparable iterations, an observed mean is about 13.3 points per iteration. A 40-point release backlog might require approximately three iterations at that historical rate, but the range widens with unknown dependencies, availability changes, defects and new scope. Treat the point values as team-local and the forecast as revisable.

\end{tcolorbox}
\section{Backlog item lifecycle and decision rights}
\begin{longtable}{|p{0.20\linewidth}|p{0.29\linewidth}|p{0.37\linewidth}|}\hline
\textbf{Stage} & \textbf{Exit evidence} & \textbf{Decision authority} \\ \hline\endfirsthead\hline
\textbf{Stage} & \textbf{Exit evidence} & \textbf{Decision authority} \\ \hline\endhead
Idea & Role, opportunity, source, broad scope & Product representative \\ \hline
Candidate & Value, assumptions, dependencies, risk & Product representative with team \\ \hline
Ready for planning & Near-term examples, size, valid slice, constraints & Team + product representative \\ \hline
In progress & Tasks owned, tests being implemented & Delivery team \\ \hline
In verification & Demonstrable increment and test evidence & Delivery team + tester \\ \hline
Accepted & Agreed outcomes met, exceptions dispositioned & Authorized product representative \\ \hline
Released & Deployment record and operational readiness & Release authority \\ \hline
\end{longtable}
\section{Nonfunctional constraints, governance and traceability}
A story may help expose a quality need, but a pervasive quality constraint should also exist as an individually traceable requirement with scope, measurable target, acceptance method and owner. Example: "For 95\% of submissions of up to 1 MB under the defined load profile, the system confirms receipt within 2 seconds." This is an instructional target, not a production SLA. Record security requirements such as least privilege, audit logging and data retention separately where necessary; link them to affected stories and tests.

For regulated or long-lived systems, link goal -> role/use case -> story -> acceptance test -> implementation -> build/test evidence -> approval. Preserve architecture decisions and compliance obligations beyond story completion. The original book distinguishes stories, use cases, scenarios and older IEEE 830-style specifications; this suite treats these as potentially complementary artifacts depending on assurance needs.

\section{Smells and corrective actions}
\begin{longtable}{|p{0.20\linewidth}|p{0.29\linewidth}|p{0.37\linewidth}|}\hline
\textbf{Symptom} & \textbf{Likely cause} & \textbf{Intervention} \\ \hline\endfirsthead\hline
\textbf{Symptom} & \textbf{Likely cause} & \textbf{Intervention} \\ \hline\endhead
A story says build an API & Technical activity confused with value & Attach task to valued outcome or track explicit enabling work \\ \hline
Story cannot be demoed independently & Layer-based decomposition & Reslice through interface, behavior and persistence \\ \hline
Acceptance is subjective & Missing test oracle & Specify concrete input, state and observation \\ \hline
Customer cannot prioritize & Unclear accountability or outcomes & Assign decision owner, clarify release goal \\ \hline
Stories repeatedly carry over & Oversize, external dependencies or unplanned work & Reduce size, investigate dependency, revisit capacity \\ \hline
Hundreds of stale details & Premature specification & Use progressively elaborated horizon and archive obsolete notes \\ \hline
Security rule appears nowhere in tests & Unmanaged cross-cutting concern & Create testable constraint and coverage links \\ \hline
\end{longtable}
\section{Implementation-ready checklist}
\begin{itemize}[leftmargin=*,itemsep=2pt]
\item A real role and valued outcome are explicit.
\item The slice is independently demonstrable or its dependency is justified.
\item Acceptance examples include relevant negative, edge and security behavior.
\item Open decisions have accountable owners.
\item Scope, risk and size are understandable to delivery participants.
\item Cross-cutting constraints are measurable and linked.
\item Test evidence and acceptance responsibility are identified.
\end{itemize}
\section{Bibliographic note}
Source: Mike Cohn, User Stories Applied: For Agile Software Development, Addison-Wesley, 2004, ISBN 0-321-20568-5. For chapter locations use the study guides and reading crosswalk. The text here is independently authored educational commentary; it does not reproduce the book’s pages.


\end{document}
